\documentclass[11pt]{article}
\usepackage{amsmath,amssymb,amsthm}
\newtheorem{theorem}{Theorem}
\title{Anti-Target Butterfly Routing: Exact Equivalence and Finite Repairs}
\author{}
\date{November 24, 2026}
\begin{document}
\maketitle

\paragraph{Status.}
The routing identity below has an elementary complete proof. The finite
feasibility statements and the one all-small-parameter repair have exact
rational certificates checked without a feasibility solver. None is
Lean formalized. No universal asymptotic decoder or global bound is
claimed, and no model definition changes.

Write $H=2^d$ and index switches by $(k,s,g)$, where $k$ is the updated
bit, $s$ is the already-updated lower prefix, and $g$ is the unchanged
higher input block. Let $T_{k,s,g}\in\{0,1\}$ be target labels.
Direct routing uses switch bits $T$; anti-target routing uses $1-T$.
In both cases the requested sign remains $1-2T_{k,s,g}$.

\begin{theorem}[Input-complement identity]
Put $\overline g=g\mathbin{\mathrm{xor}}(2^{d-k-1}-1)$ and
$S_{k,s,g}=T_{k,s,\overline g}$. Then
\[
 \pi_{1-T}(h)=\pi_S(h\mathbin{\mathrm{xor}}(H-1)).
 \tag{1}
\]
For the signed Haar query matrices on any fixed output-pole pool and
query times depending only on $(k,s)$,
\[
 A^{\rm anti}(T)_{(k,s,g),y}
 =-A^{\rm direct}(S)_{(k,s,\overline g),y}.
 \tag{2}
\]
Consequently the two families have identical feasible-pattern counts
when arbitrary signed residues are allowed, after relabeling blocks.
Positive anti-target residues correspond to negative direct residues.
\end{theorem}
\begin{proof}
Compare the anti-target run from $h$ with the $S$ run from its bitwise
complement. Before stage $k$, the lower bits agree and the remaining
bits are complementary. Their switch indices are $(k,s,g)$ and
$(k,s,\overline g)$. The first flips by $1-T_{k,s,g}$ and the second
by $S_{k,s,\overline g}=T_{k,s,g}$. Complementary incoming bits
therefore become equal. This proves the invariant inductively and
gives (1).

Input complementation sends the Haar vector $u_{k,g}$ to
$-u_{k,\overline g}$. The target signs at the corresponding relabeled
queries agree. Substituting (1) into each matrix entry gives (2).
Multiplying all residues and, if present, all block biases by $-1$
proves the feasibility assertion. Retaining queries by a condition
depending only on $(k,s)$ preserves this equivalence.
\end{proof}

\paragraph{Exact positive finite repairs.}
Use
\[
 p_y(\epsilon)=\sum_{j=0}^{d-1}y_j\epsilon^{2^j},
 \qquad t_{k,s}=1-p_s(\epsilon).
\]
At $\epsilon=1/100$, all 4096 target words for $H=8$ have strictly
positive residues realizing all 12 queries under anti-target routing,
without block biases. The saved certificate has 256 square active
bases. For every one of the 16 choices of the four bottom target bits,
the signed query matrix is exactly unchanged; the checker verifies
all of these equalities and all square-basis witnesses.

At $H=16$ the saved old direct-routing positive obstruction, and 19
additional random target words, all have positive anti-target witnesses
at the same finite parameter. Each witness satisfies all 32 signed
query scores and all 16 residue inequalities with margin at least one.
The random sample is not an estimate certified for growing $H$.

\paragraph{The old sixteen-pole word has an all-small repair.}
For the target in
\texttt{butterfly-adaptive-haar-positive-h16-dual.json}, the first
saved anti-target active basis is a square rational matrix in
$\epsilon$. Guarded exact series elimination at precision 768 gives
nonzero pivot orders
\[
 0,0,0,0,0,0,0,2,2,2,6,6,8,17,19,28.
\]
All 32 query scores and all 16 residues have positive Laurent leading
coefficients. The companion checker independently reconstructs the
finite basis, checks the guards and residuals, and compares every
leading term with the saved certificate. Invertibility over the
rational-function field and positivity of these finitely many leading
coefficients give a common punctured interval on which the repair is
strict. This is a repair of one target, not a universal sixteen-pole
or growing-order theorem.

\paragraph{Why the sign audit matters.}
A symmetric root-pair difference has a positive-definite Cauchy
kernel after reversing its routing orientation. That single-node fact
does not automatically preserve descendant signs. More generally,
(2) shows that complementing every routing bit is not a new signed
decoder family. Positive-only success or failure must be separated
from feasibility in the original model, which permits signed
numerator amplitudes.

\paragraph{A positive anti-target obstruction at thirty-two poles.}
Take the same 49 retained target labels as
\texttt{butterfly-underloaded-h32-signed.tex}, and complement every
routing bit while preserving those labels. For every sufficiently
small positive $\epsilon$, this fixed anti-target routing has no
nonnegative-residue solution, without block biases. The companion
\texttt{butterfly-antitarget-h32-obstruction-check.py} verifies a
strictly positive dependence among 31 retained signed query rows and
the two coordinate rows for $c_0\geq0$ and $c_1\geq0$.

The saved 33-column support has rank 32. Appending the normalization
that the coefficients sum to one gives a square rational system.
The checker reconstructs its positive rational solution at
$\epsilon=1/100$ and independently verifies the zero weighted sum.
Guarded precision-1536 elimination then gives 33 positive Laurent
leading coefficients, of orders between zero and 72. Exactly as for
the primal continuation above, this proves a positive dependence
on a common punctured interval. Taking its inner product with a
hypothetical nonnegative residue vector contradicts the 31 strict
query inequalities.

This excludes neither signed residues nor other auxiliary routing
completions. In particular, the direct routing for these same target
labels has a signed all-small-parameter solution, certified in
\texttt{butterfly-underloaded-h32-signed.tex}. Thus neither uniform
choice of direct or anti-target routing supplies a universal positive
pool by the root positive-definiteness argument alone.
\end{document}
